\documentclass{article}

\usepackage{amsmath,amssymb}
\usepackage{xurl}
\usepackage[hidelinks]{hyperref}

% Shared commands for notes in docs/paper-gaps/.

\DeclareUnicodeCharacter{2080}{\ifmmode{}_{0}\else\textsubscript{0}\fi}
\DeclareUnicodeCharacter{2081}{\ifmmode{}_{1}\else\textsubscript{1}\fi}
\DeclareUnicodeCharacter{2082}{\ifmmode{}_{2}\else\textsubscript{2}\fi}
\DeclareUnicodeCharacter{2099}{\ifmmode{}_{n}\else\textsubscript{n}\fi}

\newcommand{\Lean}{\textsc{Lean}}
\ExplSyntaxOn
\NewDocumentCommand{\leanid}{m}{
  \ifmmode
    \textsf{\detokenize{#1}}
  \else
    \group_begin:
    \urlstyle{sf}
    \tl_set:Nn \l_tmpa_tl {#1}
    \tl_replace_all:Nnn \l_tmpa_tl {\_} {_}
    \exp_args:NV \path \l_tmpa_tl
    \group_end:
  \fi
}
\ExplSyntaxOff
\providecommand{\tr}{\operatorname{tr}}
\newcommand{\tnleanIssueBase}{https://github.com/LionSR/TNLean/issues/}
\newcommand{\tnleanPrBase}{https://github.com/LionSR/TNLean/pull/}
\newcommand{\tnleanDocBase}{https://sirui-lu.com/TNLean/docs/}
\newcommand{\ghissue}[1]{\href{\tnleanIssueBase#1}{GitHub issue \##1}}
\newcommand{\ghpr}[1]{\href{\tnleanPrBase#1}{PR \##1}}
\newcommand{\doclink}[1]{\href{\tnleanDocBase#1.html}{\path{#1}}}

% Dirac notation and the identity operator, as in blueprint/src/macros/common.tex.
% \providecommand keeps note-local definitions authoritative.
\usepackage{dsfont}
\providecommand{\ket}[1]{|#1\rangle}
\providecommand{\bra}[1]{\langle #1|}
\providecommand{\braket}[2]{\langle #1 | #2 \rangle}
\providecommand{\ketbra}[2]{|#1\rangle\!\langle #2|}
\providecommand{\Id}{\mathds{1}}
\providecommand{\one}{\mathds{1}}
\providecommand{\C}{\mathbb{C}}
\providecommand{\MN}[1]{M_{#1}(\C)}
\providecommand{\MD}{M_D(\C)}

% External referencing for paper sources.  The build helper writes sanitized
% auxiliary files under build/paper-aux/ and excludes duplicated source labels.
\usepackage{xr}

% Import a generated paper auxiliary file with a caller-supplied label prefix.
% The candidate paths support builds from docs/paper-gaps/, the repository root,
% and build/paper-aux/ itself.
\newcommand{\importpaperaux}[2]{%
  \IfFileExists{../../build/paper-aux/#2.aux}{%
    \externaldocument[#1]{../../build/paper-aux/#2}%
  }{%
    \IfFileExists{build/paper-aux/#2.aux}{%
      \externaldocument[#1]{build/paper-aux/#2}%
    }{%
      \IfFileExists{#2.aux}{%
        \externaldocument[#1]{#2}%
      }{}%
    }%
  }%
}

% General external reference macros
\newcommand{\paperref}[2]{\ref{#1-#2}}
\newcommand{\papereqref}[2]{(\ref{#1-#2})}

% CPSV16 source (1606.00608 / MPDO-22-12-17-2.tex) external document and shortcuts
\importpaperaux{cpsv16-}{1606.00608/MPDO-22-12-17-2-tnlean}
\newcommand{\cpsvref}[1]{\paperref{cpsv16}{#1}}
\newcommand{\cpsveqref}[1]{\papereqref{cpsv16}{#1}}

% Verdict marker: \gapnote{<kind>}{<status>}. Kind is the policy
% classification (clarification, local-correction, scope-restriction,
% unfaithful, false-source, open-gap); status is open, wip (elimination
% underway), resolved, or historical. Machine-read by the site index;
% prints a verdict line.
\newcommand{\gapnote}[2]{\par\noindent\textbf{Verdict:} #1 (#2).\par}


\title{The Comparison Matrix Between the Two First-Cut Factorizations}
\author{}
\date{2026-09-07}

\begin{document}
\maketitle
\gapnote{scope-restriction}{open}

\section{The comparison step}

Let $\mathcal U$ be a matrix product operator tensor with physical dimension
$d$ and bond dimension $D$, let $M_1$ and $M_2$ be its two matrix flattenings,
and let $r$ and $\ell$ be their ranks. Section~III of
Ref.~\cite{Cirac2017MPU} constructs factorizations
$M_1=\tilde X_1\tilde Y_1$ and $M_2=X_2Y_2$ from singular value
decompositions, together with the auxiliary factors $\tilde Z_1$ and $Z_2$
that invert $\tilde Y_1$ and $Y_2$ on the right. Writing $\rho$ for the
positive diagonal fixed point of the normalized transfer map and
$W_\rho=\operatorname{Id}_d\otimes\rho$ for the associated weight, the printed
normalizations are
\begin{align}
    X_2^\dagger X_2
     & =\operatorname{Id}_\ell,
    \label{eq:fbc25_inverse_compatible_comparison_matrix_second_isometry} \\
    \tilde X_1^\dagger W_\rho\tilde X_1
     & =\operatorname{Id}_r.
    \label{eq:fbc25_inverse_compatible_comparison_matrix_weighted_isometry}
\end{align}
The second factorization is a plain isometry, whereas the first is an isometry
only against the weight $W_\rho$.\footnote{Source traceability:
    \path{Papers/1703.09188/paper_v2.tex}, lines 480--500. The declared
    auxiliary-label build for this archived source fails on the title's
    malformed accent command, so the printed section designation is used.}

The appendix of Ref.~\cite{FrancoRubio2025MPUGauging} treats a tensor whose
adjoint is a unitary virtual gauge conjugate of itself,
$\mathcal U^\dagger_{ij}=T^\dagger\mathcal U_{ij}T$, and builds a second
factorization $M_1=X_1Y_1$ out of the second-cut factors,
\begin{align}
    X_1
     & =A Y_2^\dagger,
    \label{eq:fbc25_inverse_compatible_comparison_matrix_candidate_x} \\
    Y_1
     & =X_2^\dagger B,
    \label{eq:fbc25_inverse_compatible_comparison_matrix_candidate_y}
\end{align}
where $A$ and $B$ are the virtual dressings built from $T$. It then compares
the two factorizations of $M_1$ through a single invertible matrix $K$, printed
as
\begin{align}
    K
     & =\tilde X_1^\dagger X_1
    =\tilde Y_1Y_1^\dagger,
    \label{eq:fbc25_inverse_compatible_comparison_matrix_printed_k} \\
    \tilde Y_1
     & =KY_1,
    \label{eq:fbc25_inverse_compatible_comparison_matrix_printed_pairing} \\
    \tilde X_1
     & =X_1K^\dagger.
    \label{eq:fbc25_inverse_compatible_comparison_matrix_printed_adjoint}
\end{align}
This note concerns
(\ref{eq:fbc25_inverse_compatible_comparison_matrix_printed_k}) and
(\ref{eq:fbc25_inverse_compatible_comparison_matrix_printed_adjoint}).\footnote{Source
    traceability: \path{Papers/2502.20257/main.tex}, lines 5432--5443, in the
    proof of the second appendix proposition of the $\mathbb Z_2$ appendix.}
The two deviate from the formal statement in different ways. The first printed
expression for $K$ needs the weight $W_\rho$ inserted; that is a local
correction, and the formal development applies it, so nothing about it remains
open. The second names $K^\dagger$ as the inverse of $K$, hence asserts that
$K$ is unitary, and the formal development proves no such thing: it exhibits a
separately constructed two-sided inverse instead. The formal statement is
therefore a restriction of the printed one, and that restriction, not the
weight, fixes the classification of this note.

\section{The comparison matrix is determined, and carries the weight}

Suppose $M_1=X_1Y_1=\tilde X_1\tilde Y_1$, that $\tilde L$ is a left inverse of
$\tilde X_1$, and that $R$ is a right inverse of $Y_1$. Put $K=\tilde LX_1$.
Then
\begin{align}
    \tilde Y_1
     & =\tilde LM_1
    =\tilde LX_1Y_1
    =KY_1,
    \label{eq:fbc25_inverse_compatible_comparison_matrix_derived_pairing} \\
    X_1
     & =M_1R
    =\tilde X_1\tilde Y_1R
    =\tilde X_1K,
    \label{eq:fbc25_inverse_compatible_comparison_matrix_derived_factor} \\
    K
     & =\tilde Y_1R.
    \label{eq:fbc25_inverse_compatible_comparison_matrix_derived_k}
\end{align}
The third line follows from the second because $\tilde L\tilde X_1
=\operatorname{Id}_r$. In particular $K$ does not depend on the choice of
$\tilde L$ or of $R$: it is the unique matrix with
(\ref{eq:fbc25_inverse_compatible_comparison_matrix_derived_pairing}), since
$Y_1$ has a right inverse.

Now
(\ref{eq:fbc25_inverse_compatible_comparison_matrix_second_isometry}) and
(\ref{eq:fbc25_inverse_compatible_comparison_matrix_candidate_y}) give
$Y_1Y_1^\dagger=X_2^\dagger BB^\dagger X_2=\operatorname{Id}_\ell$, because $B$
is a unitary dressing. Hence $Y_1^\dagger$ is a right inverse of $Y_1$ and the
second printed expression in
(\ref{eq:fbc25_inverse_compatible_comparison_matrix_printed_k}) is exactly
(\ref{eq:fbc25_inverse_compatible_comparison_matrix_derived_k}); it is correct
as printed.

The first printed expression is not. By
(\ref{eq:fbc25_inverse_compatible_comparison_matrix_weighted_isometry}) the
matrix $\tilde X_1^\dagger$ is a left inverse of $\tilde X_1$ only after the
weight is inserted, so the available left inverse is
$\tilde L=\tilde X_1^\dagger W_\rho$ and the comparison matrix is
\begin{align}
    K
     & =\tilde X_1^\dagger W_\rho X_1.
    \label{eq:fbc25_inverse_compatible_comparison_matrix_weighted_k}
\end{align}
Deleting $W_\rho$ from
(\ref{eq:fbc25_inverse_compatible_comparison_matrix_weighted_k}) reproduces the
printed formula, and the two agree exactly when
$\tilde X_1^\dagger(W_\rho-\operatorname{Id}_{dD})X_1=0$, hence in particular
when $\rho=\operatorname{Id}_D$. Since $\rho$ is the normalized transfer fixed
point, that equality is not available in general. The printed
$\tilde X_1^\dagger$ is therefore the same shorthand for the weighted adjoint
that Section~III of Ref.~\cite{Cirac2017MPU} uses throughout: every occurrence
of the adjoint of $\tilde X_1$ in the role of a left inverse is an occurrence
of $\tilde X_1^\dagger W_\rho$. Reading it that way makes the appendix
comparison correct with no further change, which is why the weight is a local
correction rather than a defect in the argument.

The formal development uses
(\ref{eq:fbc25_inverse_compatible_comparison_matrix_weighted_k}) as the
definition and proves
(\ref{eq:fbc25_inverse_compatible_comparison_matrix_derived_k}) as a
theorem.\footnote{Formal declarations:
    \leanid{MPOTensor.inverseCompatibleComparisonK} and
    \leanid{MPOTensor.inverseCompatibleComparison} in
    \doclink{TNLean/MPS/MPU/InverseCompatibleCutComparison}; the weighted
    normalization is \leanid{MPOTensor.sourceX₁_weighted_isometry}.}

\section{The inverse comparison}

Equation
(\ref{eq:fbc25_inverse_compatible_comparison_matrix_printed_adjoint})
asserts more than invertibility of $K$: it names $K^\dagger$ as the inverse,
which is the statement that $K$ is unitary. In the appendix that unitarity is
obtained afterwards, from the pleasant properties reestablished for the new
factorization.

The same construction applied in the opposite direction supplies an inverse
without any normalization argument. If $L$ is a left inverse of $X_1$ and
$\tilde Z_1$ a right inverse of $\tilde Y_1$, put $J=L\tilde X_1$. Then, by
(\ref{eq:fbc25_inverse_compatible_comparison_matrix_derived_pairing})--(\ref{eq:fbc25_inverse_compatible_comparison_matrix_derived_k})
with the two factorizations exchanged,
\begin{align}
    Y_1
     & =J\tilde Y_1,
    \label{eq:fbc25_inverse_compatible_comparison_matrix_reverse_pairing} \\
    \tilde X_1
     & =X_1J,
    \label{eq:fbc25_inverse_compatible_comparison_matrix_reverse_factor} \\
    J
     & =Y_1\tilde Z_1,
    \label{eq:fbc25_inverse_compatible_comparison_matrix_reverse_j}
\end{align}
and substituting each pairing into the other gives
\begin{align}
    KJ
     & =\operatorname{Id}_r,
    \label{eq:fbc25_inverse_compatible_comparison_matrix_kj} \\
    JK
     & =\operatorname{Id}_\ell.
    \label{eq:fbc25_inverse_compatible_comparison_matrix_jk}
\end{align}
Thus $J$ occupies the position of $K^\dagger$ in
(\ref{eq:fbc25_inverse_compatible_comparison_matrix_printed_adjoint}), and
$J=K^\dagger$ holds precisely when $K$ is unitary. The formal statement keeps
$J$ as constructed and asserts
(\ref{eq:fbc25_inverse_compatible_comparison_matrix_reverse_pairing})--(\ref{eq:fbc25_inverse_compatible_comparison_matrix_jk})
only; it asserts neither $J=K^\dagger$ nor unitarity of $K$. It also keeps the
rectangular index types $r$ and $\ell$ rather than transporting along the
proved equality $r=\ell$.\footnote{Formal declarations:
    \leanid{MPOTensor.inverseCompatibleComparisonJ} and
    \leanid{MPOTensor.rightRank_eq_leftRank_of_physicalAdjointTensor_eq_unitary_gauge}
    in \doclink{TNLean/MPS/MPU/InverseCompatibleCutComparison}.}

\section{Verdict}

The second printed expression for the comparison matrix is correct. The first
is the usual shorthand in which the adjoint of $\tilde X_1$ stands for the
weighted adjoint $\tilde X_1^\dagger W_\rho$, and it is correct only under that
reading; the formalization inserts the weight. This is a local correction, it
changes no conclusion of the appendix argument, and it is carried out in full,
so it leaves nothing open.

The identification of the inverse comparison with $K^\dagger$ is a separate
statement that the appendix obtains from the pleasant properties of the new
factorization. The formal development proves only that $K$ has a two-sided
inverse, exhibited explicitly. Its pairing is therefore the printed pairing
restricted to the case in which the inverse is that explicit matrix rather than
$K^\dagger$, and the unitarity of $K$ that the appendix establishes afterwards
is absent. This restriction, and not the weight, is the live content of the
note, and it is the reason the note is classified as a scope restriction rather
than as a local correction. The elimination step is the formalization of the
appendix normalization that makes $K$ unitary; the note becomes resolved when
that argument gives $J=K^\dagger$, and until then the $K^\dagger$ form of the
pairing remains outside the formalized statement and the note stays open.

\bibliographystyle{alpha}
\bibliography{references}

\end{document}
